\vspace{8pt}
{\scriptsize\setlength{\parskip}{2pt}\textbf{Références.}
André et Meslin (2023), \emph{Insee Analyses} n°~91.
André et Sireyjol (2021), « Effets redistributifs de l'imposition des couples et des familles : une étude par
microsimulation de l'impôt sur le revenu », \emph{Économie et Statistique} n°~526-527.
Assemblée nationale (2026), commission d'enquête sur l'imposition des plus hauts patrimoines et des revenus les
plus élevés, rapport n°~3056.
Avenancio-León et Howard (2022), « The Assessment Gap », \emph{QJE} 137(3).
Bach, Bozio, Fabre, Guillouzouic, Leroy et Malgouyres (2019), note IPP n°~46.
Bach, Bozio, Grimprel, Guillouzouic, Landais et Malgouyres (2025), Focus du CAE n°~118.
Bachelet et Poulhès (2017), CGDD.
Banzhaf et Lavery (2010), \emph{Journal of Urban Economics} 67(2).
Beer, de Mooij et Liu (2020), \emph{Journal of Economic Surveys} 34(3).
Bérard et Trannoy (2018), \emph{Économie et Statistique} n°~500-501-502.
Bernini et Treibich (2016), \emph{Industrial and Corporate Change} 25(5).
Berry (2021), « Reassessing the Property Tax », Harris School.
Best et Kleven (2018), \emph{Review of Economic Studies} 85(1).
Bontemps, Cherbonnier et Magnac (2023), TSE WP 23-1486.
Cabannes, Houdré et Landais (2014), \emph{Économie et Statistique} 467-468.
CAE (2025, 2026), Focus n°~124 et n°~147.
Carbonnier (2021), « Imposition jointe des revenus et emploi des femmes mariées : estimation à partir du cas
français », \emph{Revue économique} 72(2).
Chetty (2009), « Is the Taxable Income Elasticity Sufficient to Calculate Deadweight Loss? », \emph{AEJ:
Economic Policy} 1(2).
Chetty, Friedman, Olsen et Pistaferri (2011), \emph{QJE} 126(2).
Chetty, Guren, Manoli et Weber (2013), « Does Indivisible Labor Explain the Difference between Micro and Macro
Elasticities? », \emph{NBER Macroeconomics Annual} 27.
Chodorow-Reich, Smith, Zidar et Zwick (2026), NBER WP 32180 (2024), version de janvier 2026.
Conseil constitutionnel, décisions 2000-437 DC, 2012-662 DC et 2019-796 DC.
Conseil d'État, 13 mars 2025, n°~481538.
Conseil des prélèvements obligatoires (2017, 2023, 2025), « Adapter l'impôt sur les sociétés à une économie
ouverte » ; « Les différences d'imposition sur les bénéfices entre PME et grandes entreprises » ; « Tracer un
cadre fiscal et social pluriannuel pour l'industrie française ».
Dachis, Duranton et Turner (2012), \emph{Journal of Economic Geography} 12(2).
Dolls, Fuest, Krolage et Neumeier (2025), \emph{Journal of Urban Economics} 145.
Edmark, Liang, Mörk et Selin (2016), « The Swedish Earned Income Tax Credit: Did It Increase Employment? »,
\emph{FinanzArchiv} 72(4).
Eerola, Harjunen, Lyytikäinen et Saarimaa (2021), \emph{Journal of Urban Economics} 124.
Feldstein (1999), \emph{REStat} 81(4).
Ferreira (2010), \emph{Journal of Public Economics} 94(9-10).
France Stratégie (2023), comité d'évaluation des réformes de la fiscalité du capital, quatrième rapport.
Fritzsche et Vandrei (2019), \emph{Regional Science and Urban Economics} 74.
Harberger (1964), « Taxation, Resource Allocation, and Welfare », dans NBER et Brookings Institution, \emph{The
Role of Direct and Indirect Taxes in the Federal Revenue System}, Princeton University Press.
Hines (1999), « Three Sides of Harberger Triangles », \emph{Journal of Economic Perspectives} 13(2).
HMRC (2023), évaluation de l'allègement pour primo-accédants.
Høj, Jørgensen et Schou (2018), \emph{Fiscal Studies} 39(2).
Hourani, Millar-Powell, Perret et Ramm (2023), OCDE, \emph{Taxation Working Paper} n°~65.
IFS (2023), \emph{Green Budget}, ch.~10.
IPP (2023), rapport n°~47.
Kleven et Kreiner (2006), « The Marginal Cost of Public Funds: Hours of Work versus Labor Force
Participation », \emph{Journal of Public Economics} 90(10-11).
Koster et Pinchbeck (2022), \emph{AEJ: Economic Policy} 14(1).
Lefebvre, Lehmann et Sicsic (2025), \emph{Scandinavian Journal of Economics} 127(2) (IZA DP 16112, 2023).
Lehmann, Marical et Rioux (2013), « Labor Income Responds Differently to Income-Tax and Payroll-Tax Reforms »,
\emph{Journal of Public Economics} 99.
Lutz (2015), \emph{AEJ: Economic Policy} 7(1).
Määttänen et Terviö (2022), \emph{Economic Journal} 132(644).
Maffini, Xing et Devereux (2019), \emph{AEJ: Economic Policy} 11(3).
Martínez, Saez et Siegenthaler (2021), \emph{AER} 111(2).
Morin (2014), \emph{Insee Première} 1492.
Nielsson, Wroblewski et Yding (2024), document de travail.
OCDE, \emph{Revenue Statistics}, \emph{Corporate Tax Statistics} (2025), \emph{Taxing Wages} (2025).
Paquier et Sicsic (2022), \emph{Économie et Statistique} n°~530-31.
Sénat (1977), M.~Blin, rapport général sur le PLF pour 1978, n°~88 ; Sénat (2026), mission d'information sur le
poids des prélèvements obligatoires en France et ses conséquences sur la compétitivité des entreprises, rapport
n°~942 (2025-2026).
Sicsic (2022), « Does Labour Income React More to Income Tax or Means-Tested Benefits Reforms? », \emph{Fiscal
Studies} 43(3).
Tax Foundation, \emph{International Tax Competitiveness Index} 2025 et 2026 ; \emph{Dividend Tax Rates in
Europe} 2026.
Yagan (2015), \emph{AER} 105(12).
Zwick et Mahon (2017), \emph{AER} 107(1).
Textes : LF 2024 (art.~19 et 33), LF 2025, LFSS et LF 2026, PLF 2027 (art.~5).
Détail des sources et statut de lecture : \texttt{ir\_csg/}, \texttt{foncier/notes/}, \texttt{is/notes/},
\texttt{pfu/notes/}, \texttt{revue/}.\par}
